% !TEX program = xelatex

\documentclass[11pt,a4paper]{ctexart}

% ====== 字体 ======
\usepackage{fontspec}
\setmainfont{Latin Modern Roman}
\usepackage{unicode-math}
\setmathfont{Latin Modern Math}

% ====== 页面设置 ======
\usepackage[margin=2.5cm]{geometry}
\usepackage{setspace}
\onehalfspacing
\pagestyle{plain}

% ====== 数学宏包（三个经典功能，但用 unicode-math 替代了 amssymb） ======
\usepackage{mathtools}          % 必须放在 unicode-math 之前
\usepackage{amsmath}            % 已被 mathtools 自动载入，保留显式声明无妨
% \usepackage{amssymb}         % 删除！与 unicode-math 冲突，且 unicode-math 已提供所有符号

% ====== 三线表 + 数字对齐 ======
\usepackage{booktabs}
\usepackage{siunitx}
\usepackage{enumitem}          % 增强列表控制

% ====== 交叉引用与超链接 ======
\usepackage[colorlinks=true,
            linkcolor=blue,
            citecolor=green,
            urlcolor=cyan]{hyperref}
\usepackage[noabbrev,nameinlink]{cleveref}  % 必须最后加载

% ====== 自定义命令 ======
\newcommand{\dd}{\mathrm{d}}
\newcommand{\source}[1]{\par \hfill ——#1}

\title{红色文化与科技报国交织的宁河底色}
\author{王晨暄}
\date{\today}

\begin{document}
\maketitle

宁河是一方红色热土，更是一方科技报国的精神高地。于方舟烈士创建天津党团组织，组织“新生社”
与周恩来领导的“觉悟社”并肩传播革命思想，被俘后宁死不屈，年仅28岁英勇就义，
彰显了宁河青年心怀家国、挺膺担当的品格。河西大会战中，宁河举全县之力改土治碱、兴修水利，
建成了华北平原上的“鱼米之乡”，展现了迎难而上、艰苦奋斗的集体精神。两种精神在此交汇——
“不畏艰苦，勇于担当；齐心奋斗，百折不挠”的宁河血脉，既熔铸于革命年代的慷慨赴义，
也赓续于科技报国的使命传承之中。

而于敏，正是这条精神血脉在科技报国征程上的集中体现。“愿将一生献宏谋”——从宁河芦台镇走出的
“两弹一星”元勋，隐姓埋名二十八载，带领团队从零起步，攻克氢弹理论难题。于敏构型的横空出世，
使中国以世界惊叹的速度实现了从原子弹到氢弹的跨越。但他一生淡泊名利，三次婉拒“氢弹之父”的称谓，
将荣誉归于集体。“身为一叶无轻重”，于敏用毕生践行了这句自勉之言。他身上所凝聚的，
正是宁河血脉中“心怀家国、默默奉献、勇攀高峰”的精神品格。这份精神并未止步于他本人，
在他的母校芦台一小，“于敏纪念室”与“于敏科学乐园”已成为播撒科学种子、传承报国情怀的重要阵地，
让“胸怀家国、艰苦奋斗”的品质在年轻一代心中生根发芽、代代赓续。

于敏所代表的科技报国精神，在当代宁河接续燃烧。新天钢联合特钢有限公司以“实业立国、复兴中华”为初心，
积极探索绿色高效的生产模式，攻坚高品质钢材锻造，在传承艰苦奋斗精神的同时勇于超越。天津宁河
低空产业园则集合二十余家企业、联手高校，在新型材料、零部件等领域不断突破，
探索低空经济新路径，书写科技兴国、科技造福民生的新篇章。

从革命烈士“苟利国家生死以，岂因祸福避趋之”的慷慨赴义，到科研工作者“誓干惊天动地事，甘做隐姓埋名人”
的无私奉献；从芦台一小“孝悌忠信，勤俭是纲”的薪火相传，到新时代“科技兴国、实业报国”的使命赓续——
变的是时代课题，不变的是宁河人迎难而上的精神品格与科技报国的赤子之心。红色文化与科技报国于此交织，
熔铸成“心怀家国，踔厉奋发”的宁河精神底色。

\end{document}